\section{Annex}

\subsection{Survey Design (Direct Assessment, Idiom Replacement)}
\label{appendix:survey-design}

This appendix reproduces the full survey instrument described in Chapter~\ref{ch:methodology} (protocol) and Section~\ref{subsec:rq3-human-eval} (instantiation). It documents the per-item layout, per-respondent flow, quality-control item construction, fixed-threshold filter values, and the form-deployment choice.

\subsubsection{Per-item layout}

Each rated item presents the original sentence, the system-detected idiomatic expression, and the system's full-sentence replacement, followed by three independent 0--100 sliders. An optional free-text field captures alternative paraphrases for qualitative analysis.

\begin{verbatim}
Original sentence:
  "Unni, the stylist, is on cloud nine after having an opportunity to
   style the beard of his favourite star."

Idiomatic expression detected by the system:
  "on cloud nine"

System's replacement of the full sentence:
  "Unni, the stylist, is extremely happy after having an opportunity to
   style the beard of his favourite star."

Please rate the system's replacement on three dimensions
(0 = unacceptable, 100 = perfect):

  Grammaticality (is it correct English?)            [ 0 -- 100 ]
  Meaning preservation (same meaning as original?)   [ 0 -- 100 ]
  Simplicity (clearer than the idiom version?)       [ 0 -- 100 ]

Optional --- if the replacement misses the mark,
what would you write instead?

  [                                                                  ]
\end{verbatim}

The detected expression is shown explicitly so annotators do not have to find the idiom themselves; this controls for ``did the annotator notice the idiom?'' as a confound. The SemEval Task B reference paraphrase is \textit{not} shown: Task B references are comparison sentences, not meaning-equivalent paraphrases, and showing them would distort the meaning-preservation rating.

\subsubsection{Per-respondent flow}

\begin{enumerate}
    \item \textbf{Intro page} --- plain-language task description, ethics paragraph (supplied by Mari Carmen Su\'arez-Figueroa), and a native-English-speaker self-attestation. No name or email is collected; the only demographic field is the binary native-English-speaker check.
    \item \textbf{Calibration item} --- one worked example before the rated items, with explanations of each rubric dimension. Not graded; for orientation only.
    \item \textbf{Rated items (16)} --- random subset of the 90-item pool ($30~\text{sources} \times 3~\text{conditions}$). Sampling is constrained so each respondent sees a mix of the three system conditions and never the same source twice except as a deliberate repeat-pair (Section~\ref{appendix:qc-construction}).
    \item \textbf{QC items (4)} --- two bad-reference items and two repeat-pair items, interleaved among the 16 rated items at randomised positions. The total of 20 items per respondent matches the 20\% QC ratio of \cite{Graham2013}.
    \item \textbf{Submission} --- on submit, raw scores are written to a Google Sheet for offline analysis. No real-time feedback is shown to the respondent.
\end{enumerate}

Estimated completion time is 22--25 minutes. A two-respondent pilot is run before live distribution to verify median completion time; if median exceeds 25 minutes the rated-item count is cut to 12 (with the QC count held at 4) before going live.

\subsubsection{Quality-control item construction}
\label{appendix:qc-construction}

\paragraph{Bad-reference items.} Two of each respondent's 20 items are constructed by deliberate degradation of a system output, anchored on \textit{meaning preservation} as the most binary and least ambiguous of the three rubric dimensions. The degradation strategies, in order of preference:
\begin{enumerate}
    \item Replace the whole sentence with a sentence whose meaning differs from the original (default).
    \item Replace the idiom with a meaningless or syntactically broken phrase (breaks grammaticality instead of meaning).
    \item Replace the idiom with a more obscure idiom (breaks simplicity instead of meaning).
\end{enumerate}
Bad references are generated by a separate LLM run with a ``make this output worse on dimension X'' prompt and manually reviewed by the author before survey distribution. Each respondent sees one bad reference of strategy~1 and one of either strategy~2 or~3, balanced across the respondent pool.

\paragraph{Repeat-pair items.} Two items in each respondent's set are duplicates of two other items in the same set, with at least five other items separating each pair's two showings. The pairs are randomised per-respondent so the same source items are not always re-rated, avoiding bias in the candidate pool.

\subsubsection{Fixed-threshold filter values}

In place of the per-worker difference-of-means t-test of \cite{Graham2017} (which has effectively no power at $n = 2$ per QC type, see Section~\ref{sec:qc-filter-design}), two fixed-threshold filters are applied. Threshold values are fixed in advance and reported here:

\begin{itemize}
    \item \textbf{Bad-reference filter.} Exclude a respondent if \textit{both} of their bad-reference items score $\geq$ the corresponding system output on the targeted rubric dimension. Allowing one miss accommodates accidental misclicks; two misses is signal that the respondent is not discriminating.
    \item \textbf{Repeat-pair filter.} For each repeat pair, compute the absolute score difference per rubric dimension. Exclude a respondent whose mean repeat-pair difference, averaged across the two pairs and the three rubric dimensions, exceeds 25 points (out of 100).
\end{itemize}

\subsubsection{Distribution}

The survey is deployed as a Google Form with an Apps Script back-end that implements per-respondent randomised item-subset selection and QC injection. Raw responses dump to a Google Sheet which is exported as CSV for offline aggregation. Distribution is via Mari Carmen Su\'arez-Figueroa's network of native-English-speaker contacts within the UPM Easy-to-Read research group's ecosystem.

The instrument uses 0--100 numeric input fields rather than continuous-slider widgets (Google Forms does not support a true continuous slider). Per Graham et al.\ \cite{Graham2013}, the load-bearing property of the protocol is the \textit{continuity} of the rating scale rather than its render-time resolution; integer 0--100 input is faithful to the protocol.

\newpage
